\documentclass[UTF8,a4paper,11pt,fontset=fandol]{ctexart}

\usepackage{geometry}
\geometry{top=2.4cm,bottom=2.4cm,left=2.6cm,right=2.6cm}
\usepackage{booktabs}
\usepackage{tabularx}
\usepackage{longtable}
\usepackage{array}
\usepackage{enumitem}
\usepackage{amsmath}
\usepackage{float}
\usepackage{xcolor}
\usepackage{hyperref}
\hypersetup{colorlinks=true,linkcolor=black,urlcolor=blue}
\usepackage{titlesec}
\titleformat{\section}{\Large\bfseries}{\thesection}{0.8em}{}
\titleformat{\subsection}{\large\bfseries}{\thesubsection}{0.7em}{}
\usepackage{fancyhdr}
\pagestyle{fancy}
\fancyhf{}
\cfoot{\thepage}
\renewcommand{\headrulewidth}{0pt}
\setlength{\parindent}{2em}
\setlength{\parskip}{0.3em}

\newcommand{\code}[1]{\texttt{\detokenize{#1}}}

\begin{document}

\begin{center}
  {\zihao{2}\bfseries 开题报告材料说明\par}
  \vspace{0.4em}
  {\small 2026 年 11 月 · 供后续接手写 / 改开题报告的人使用\par}
\end{center}

\vspace{0.5em}

给后面接手写 / 改开题报告的人看的。先读这一份，再按需打开 \code{开题报告.pdf} 改内容，不用从头翻代码和实验记录。

当前状态（2026-11-25）：软件闭环基本完成，识别与拒识都有实测数字；真人命令正样本、硬件实机、演示视频和成本表还没做。开题报告里这些要按“待完成”写，不要写成已经完成。

\section{这份材料里有什么}

\begin{table}[H]
  \centering
  \small
  \begin{tabularx}{\textwidth}{@{}lX@{}}
    \toprule
    文件 & 用途 \\
    \midrule
    \code{开题报告.tex} / \code{开题报告.pdf} & 正文（八节，封面、目录、表格、公式），直接改 \code{.tex} 后编译 \\
    \code{开题报告材料说明.tex} / \code{.pdf} & 本说明：结构对应、数字口径、旧结论黑名单、缺口、参考文献 \\
    \code{build.ps1} & 一键编译两份 PDF（xelatex 跑两遍） \\
    \code{../engineering/docs/} & 原始权威文档（测试报告、v5 方案、建议书、旧开题 PPT/PDF） \\
    \bottomrule
  \end{tabularx}
\end{table}

\section{开题报告按这个结构写}

正文已经是这个结构，可以直接用；换学校模板时按下面的对应关系搬。

\begin{table}[H]
  \centering
  \scriptsize
  \begin{tabularx}{\textwidth}{@{}p{3.2cm}Xp{4.2cm}@{}}
    \toprule
    章节 & 写什么 & 主要素材 \\
    \midrule
    一、选题背景与意义 & 方言用户（尤其老人）用不了普通话智能家居；封闭音箱难改；本装置不改固件、不依赖厂商 API，端侧识别后播普通话指令 & \code{engineering/docs/项目建议书_正式版.md} 第一节 \\
    二、研究现状与前期基础 & MFCC/DTW $\rightarrow$ 端侧 1D-CNN $\rightarrow$ 选择性分类/拒识；粤语资源少；前期语料、特征、多条路线、拒识与执行链路 & \code{research/docs/v2_report.md}、\code{v3_report.md} \\
    三、研究目标与内容 & 六条固定命令；端侧实时识别与安全拒识；普通话播报与双执行通道；真人数据采集 & 正文第三节、\code{engineering/PROGRESS.md} \\
    四、技术路线与方案 & VAD $\rightarrow$ log-mel40+$\Delta$ $\rightarrow$ 1D-CNN 3 种子 $\rightarrow$ 温度标定 $\rightarrow$ 双信号拒识 $\rightarrow$ 语义层映射 $\rightarrow$ 播报/串口 & \code{engineering/docs/v5方案设计.md}、\code{测试报告.md} \\
    五、创新点 & 识别正确性与执行安全性分开建模；训练/部署同源；小模型 + kNN 参照池的 OOD 信号；两条硬件路线共享语义口径 & 正文第五节、\code{engineering/docs/项目审查_2026-11-13.md} \\
    六、进度安排 & 8 周计划；软件侧已完成，硬件/真人/演示待做 & 正文第六节 \\
    七、成果与限制 & 原型、数据、运行时、演示；限制见正文第七节 & 正文第七节、\code{engineering/PROGRESS.md} \\
    八、参考文献 & FLEURS、MFCC、端侧 KWS、选择性分类与标定 & 本说明第八节、\code{参考文献与数据来源} 相关条目 \\
    \bottomrule
  \end{tabularx}
\end{table}

\section{现在能用的数字（2026-11-25 口径）}

写正文时优先用这一组；括号里是必须一起写的口径。

\begin{table}[H]
  \centering
  \small
  \begin{tabularx}{\textwidth}{@{}lXX@{}}
    \toprule
    指标 & 数值 & 口径 \\
    \midrule
    命令识别 test & 0.9110（483 条） & 合成 edge-tts 粤语；同音色、未见表达，不是跨说话人 \\
    命令识别 val & 0.9523（440 条） & 合成粤语 \\
    总体误接受率 FAR & 0.0702 & 真人非命令负样本（FLEURS），样本外 \\
    合成已知命令覆盖率 & 0.8205 & 合成侧代理，不是真人命令可用性 \\
    ALL\_OFF 危险误动作率 & 0.0326 & 真人负样本；目标 $\le 0.02$，未达成 \\
    动作延迟 P95 & 290.5 ms & 当前部署口径，本机 CPU，合成语料，未实机 \\
    预触发 CPU P95 / 对话完成 P95 & 42.8 ms / 1090.5 ms & 同上 \\
    模型体积 / 纯前向 & 1.35 MB / 13--16 ms 每条 & numpy 推理 \\
    C 端内核对拍 & 482/483 & 与 numpy 决策一致 \\
    软件全量测试 & 350 passed & 最近一次；11-19 时点曾为 347 \\
    会话成功率 & 0.7702 & 目标 0.85，未达成 \\
    真人命令正样本 & 0 & 必须先写明 \\
    \bottomrule
  \end{tabularx}
\end{table}

双信号门限（要写就成对写）：$\text{maxprob} \ge 0.51863$ 且 $\text{ood} \ge -0.026796$，500 条 f16 参照池、$k=5$；高风险 ALL\_OFF 另加 $\text{ood} \ge -0.019$。

\section{两条必须写清楚的限制}

\begin{enumerate}[leftmargin=2em]
  \item \textbf{识别准确率全是合成 TTS}。测试集是同音色、未见表达的 33 个文本，不是真人、不是跨说话人；真人命令识别率目前没有数据。
  \item \textbf{真实域 FAR 是“真人非命令被误接受”}，只说明安全性的一半，不能写成真人命令识别率；覆盖率 0.8205 是合成侧代理。
\end{enumerate}

\section{不要写进正文的旧结论}

\begin{itemize}[leftmargin=2em]
  \item \textbf{“方言到普通话特征映射”主线}：特征级 RMSE 0.025、余弦 0.999，但映射后命令识别 0.1667，等于六类随机，已否决。只可作为“为什么改走语义层”的对照。
  \item 旧模型/旧阈值：\code{v3_cnn_ensemble_vad}、argmax 0.8799、阈值 0.309193 / -0.022306、树莓派主控预算、142.7 ms 的旧延迟口径，都只在讲演进时当历史对照，并标注旧口径。
  \item 二期“扩充数据不提升”是真实的负结果：可以作为研究现状/前期基础，说明“不再堆数据，转去做模型与拒识”的原因。
  \item 远程仓库的 \code{pc-dialect-engine}、\code{esp32-dialect-voice} 是两个独立尝试分支，不算本项目主线，不要混进正文。
  \item 不要把 \code{szpi/} 的板载路线写成已完成：目前只有 BSP 和硬件自检编译通过，录音/上传/UI/后端都还没实现。
\end{itemize}

\section{还缺什么（按实际状态写待完成）}

\begin{enumerate}[leftmargin=2em]
  \item 真人命令采集：0 条；\code{engineering/scripts/collect_human.py} 已备，等人录。
  \item 硬件实机联调：串口 / GPIO / 舵机 / 麦克风未接；\code{szpi/} 板卡也还没上真板。
  \item 3 分钟演示视频、采购成本表、部分答辩材料。
  \item 端侧是否上 S3 没定：C 内核主机约 107 ms/条，S3 估 0.5--1 s，可能超 300 ms 预算。
  \item 快速端点切早（13/480）和三条纯镜像纠错句（剩余 43 错中 32 条）尚未根治。
\end{enumerate}

\section{参考文献与数据来源}

\subsection*{论文 / 公开资料}

\begin{enumerate}[leftmargin=2em]
  \item Davis S, Mermelstein P. Comparison of parametric representations for monosyllabic word recognition in continuously spoken sentences. \emph{IEEE Transactions on Acoustics, Speech, and Signal Processing}, 1980, 28(4): 357--366.
  \item Sainath T N, Parada C. Convolutional neural networks for small-footprint keyword spotting. \emph{Interspeech 2015}: 1478--1482.
  \item Conneau A, Ma M, Khanuja S, et al. FLEURS: Few-shot Learning Evaluation of Universal Representations of Speech. \emph{arXiv:2205.12446}, 2022.
  \item Geifman Y, El-Yaniv R. Selective Classification for Deep Neural Networks. \emph{NeurIPS 2017}.
  \item Guo C, Pleiss G, Sun Y, Weinberger K Q. On Calibration of Modern Neural Networks. \emph{ICML 2017}.
  \item Sun Y, Ming Y, Zhu X, Li Y. Out-of-Distribution Detection with Deep Nearest Neighbors. \emph{ICML 2022}.
\end{enumerate}

\subsection*{数据与许可}

\begin{itemize}[leftmargin=2em]
  \item \textbf{FLEURS}（\code{yue_hant_hk}，CC BY 4.0）：项目只用它的 test split 做真人非命令负样本，仓库没有分发音频，只保留派生的 \code{research/data/real_domain_testset/manifest.csv} 和 \code{provenance.json}。
  \item \textbf{合成粤语命令音频}：edge-tts（\code{zh-HK-HiuGaaiNeural} / \code{HiuMaanNeural} / \code{WanLungNeural}）与 canto-tts 生成，不是真人录音；文本仅做规则审查（\code{RULE_REVIEWED}）。
  \item \textbf{普通话播报音频}：edge-tts 合成的 6 条预录音频（\code{engineering/data/audio/mandarin/}）。
  \item 仓库代码和文档目前未声明开源许可证，默认保留所有权利。
\end{itemize}

\section{原始材料在哪}

\begin{itemize}[leftmargin=2em]
  \item 工程侧：\code{../engineering/PROGRESS.md}（最新进度与坑）、\code{docs/测试报告.md}（数字）、\code{docs/v5方案设计.md}（路线与阶段）、\code{docs/技术路线修订说明.md}（否决映射路线的依据）、\code{docs/交付物清单.md}（逐项状态）。
  \item 研究侧：\code{../research/docs/v3_report.md}（第三阶段完整报告）、\code{v2_report.md}（负结果）、\code{data_v3.md} / \code{data_v3_1.md}（语料与审计）。
  \item 板卡侧：\code{../szpi/README.md}、\code{../szpi/PROGRESS.md}。
  \item 旧开题材料：\code{../engineering/docs/开题报告_智能音箱外接版_Beamer.pdf} / \code{.pptx} / \code{.tex}（结构与公式可参考，数字是旧口径，不要直接抄）。
\end{itemize}

写完后如果数字有更新，改正文和本说明里的口径表，不要把新数字只留在正文里。

\end{document}
